\chapter{Referencial Teórico}

Avaliar um serviço público digital exige olhar para mais do que a tela final de
confirmação. A experiência do cidadão começa quando ele tenta entender onde está
o serviço, continua nos caminhos que percorre e termina apenas quando a demanda
é resolvida ou abandonada. Por isso, este capítulo reúne os fundamentos que
sustentam a proposta do avaliador: conceitos de serviço público, satisfação,
qualidade percebida, usabilidade, modelos brasileiros de avaliação e sinais de
dificuldade observáveis durante a navegação.

O capítulo também parte de uma tradição anterior à digitalização dos serviços:
os modelos clássicos de qualidade em serviços. Essa base é importante porque
mostra que a percepção do usuário sempre dependeu da comparação entre o que ele
esperava receber e o que efetivamente experimentou. No ambiente público digital,
essa comparação passa a envolver, além do resultado do serviço, a clareza da
informação, a confiança no canal, a facilidade de navegação, o esforço exigido e
a capacidade do órgão de responder ao problema identificado
\cite{gronroos1984service,parasuraman1988servqual,deloneMclean2003}.

\section{Serviços e Serviços Públicos}

Serviços diferenciam-se de produtos porque são mais intangíveis, ocorrem como
séries de atividades, envolvem algum grau de interação com o usuário e, em muitos
casos, são produzidos e consumidos simultaneamente. Essa característica torna a
avaliação de serviços mais dependente da experiência vivida pelo usuário, e não
apenas de propriedades observáveis de um bem material
\cite{gronroos1984service,limaPereira2020modulo1}.

No setor público, o serviço é parte do relacionamento entre cidadão e Estado. Ele
pode envolver uma entrega individualizada, como a emissão de um documento, ou uma
entrega de interesse coletivo, como segurança, educação, saúde e infraestrutura.
Por isso, a avaliação de serviços públicos possui uma dimensão adicional: além de
medir satisfação ou qualidade, ela se conecta ao exercício da cidadania, à
prestação de contas e ao aperfeiçoamento da administração pública
\cite{brasil2017lei13460,brasil2021lei14129,limaPereira2020modulo1}.

Com a digitalização, muitos serviços passaram a ser acessados por portais,
aplicativos e sistemas integrados. Essa mudança amplia o alcance da prestação
pública, mas também cria novas exigências de usabilidade, clareza,
confiabilidade e disponibilidade. A usabilidade pode ser entendida como resultado
de uso associado à efetividade, eficiência e satisfação em um contexto definido
\cite{iso9241112018}. O usuário não avalia apenas o resultado final do serviço;
ele também avalia o caminho percorrido, a linguagem utilizada, o esforço exigido,
a facilidade de navegação e a capacidade do canal digital de resolver sua
demanda.

\section{Avaliação de Serviços}

A avaliação de serviços pode ser entendida como uma investigação qualitativa ou
quantitativa sobre a qualidade do serviço e a satisfação do usuário, considerando
a relação entre expectativa e experiência concreta \cite{limaPereira2020modulo1}.
Em outras palavras, avaliar um serviço significa observar se a entrega realizada
corresponde ao que o usuário precisava, esperava ou considerava adequado.
Modelos clássicos de qualidade de serviço indicam que a percepção do usuário é
formada por lacunas entre o serviço esperado e o serviço recebido. O modelo de
qualidade percebida de \citeonline{gronroos1984service} e o modelo SERVQUAL de
\citeonline{parasuraman1988servqual} consolidaram a ideia de que qualidade em
serviços depende de fatores como confiabilidade, responsividade, segurança,
empatia e aspectos tangíveis. Embora esses modelos tenham origem no setor
privado, eles influenciaram diversas iniciativas posteriores de avaliação de
serviços públicos e serviços digitais.

No SERVQUAL, essas dimensões aparecem como tangibilidade, confiabilidade,
responsividade, garantia e empatia. A contribuição central do modelo é tratar a
qualidade percebida como diferença entre expectativa prévia e percepção após a
experiência. Para serviços públicos, entretanto, essa leitura exige adaptação: o
cidadão geralmente não escolhe entre prestadores concorrentes e muitas vezes
depende daquele serviço para exercer um direito, cumprir uma obrigação ou acessar
um benefício. Assim, a expectativa não é formada apenas por comparação de
mercado, mas também por experiências anteriores com o Estado, comunicação
institucional, relatos de outros usuários e, mais recentemente, pela comparação
com serviços digitais privados que oferecem interfaces mais rápidas e
responsivas \cite{parasuraman1988servqual,limaPereira2020modulo2}.

Nos serviços digitais, a avaliação passou a incorporar fatores ligados ao uso de
sistemas de informação, como qualidade do sistema, qualidade da informação,
qualidade do serviço, intenção de uso, satisfação do usuário e benefícios
percebidos. Os materiais da Enap destacam que, à medida que cidadãos passaram a
interagir com sistemas governamentais digitais, tornou-se necessário adotar
modelos capazes de avaliar não apenas páginas ou portais, mas a experiência do
usuário ao longo da prestação do serviço
\cite{limaPereira2020modulo2,deloneMclean2003,pedrosa2019qualidadeGovbr}.

\section{Satisfação e Qualidade Percebida}

Satisfação e qualidade percebida são conceitos próximos, mas não idênticos. A
satisfação costuma expressar um julgamento global do usuário após consumir o
serviço, enquanto a qualidade percebida pode ser decomposta em atributos mais
específicos da experiência, como rapidez, clareza das informações, atendimento,
esforço, confiabilidade e facilidade de uso. Para \citeonline{oliver2014satisfaction},
a satisfação está relacionada ao julgamento do usuário sobre o grau em que uma
experiência atendeu ou não suas expectativas.

No contexto público, a distinção entre satisfação e qualidade percebida é
importante porque uma nota geral, sozinha, nem sempre indica o que deve ser
melhorado. Um usuário pode avaliar mal um serviço por demora, linguagem confusa,
instabilidade do sistema, falta de atendimento ou excesso de etapas. Portanto,
um avaliador de serviços públicos precisa combinar uma avaliação simples, como
nota ou estrelas, com mecanismos que ajudem a qualificar o motivo da avaliação.

Essa lógica aparece na evolução dos instrumentos brasileiros de avaliação. O
BRASP considera a qualidade percebida como um elemento central para compreender
a experiência do usuário e apoiar o aprimoramento dos serviços
\cite{silva2019brasp}. A proposta de modelo unificado busca integrar qualidade
técnica, qualidade percebida, efetividade e satisfação, reduzindo o custo de
resposta para o cidadão e tornando os resultados mais úteis para a gestão
\cite{silva2022modeloUnificado}.

\section{Avaliação de Serviços Públicos no Brasil}

A avaliação de serviços públicos no Brasil possui base normativa e base técnica.
Do ponto de vista normativo, a Lei n\textsuperscript{o} 13.460/2017 estabelece
direitos dos usuários e prevê que os resultados de avaliação subsidiem a
reorientação dos serviços \cite{brasil2017lei13460}. O Decreto
n\textsuperscript{o} 9.094/2017 institui a Carta de Serviços ao Usuário e reforça
o uso de ferramenta de pesquisa de satisfação no governo federal
\cite{brasil2017decreto9094}. A Lei n\textsuperscript{o} 14.129/2021 relaciona o
governo digital à eficiência pública, à desburocratização e à participação do
cidadão \cite{brasil2021lei14129}.

A partir desse conjunto normativo, a avaliação deixa de ser apenas uma prática
gerencial desejável e passa a compor o próprio modo de prestação dos serviços
públicos digitais. A Lei de Acesso à Informação reforça a transparência como
base para acompanhamento social; a Lei n\textsuperscript{o} 13.460/2017 vincula
avaliação à proteção do usuário; a Lei do Governo Digital incorpora canais
digitais de participação; e normas infralegais detalham instrumentos de
simplificação, monitoramento e avaliação de satisfação. A Tabela
\ref{tab:marcos-legais-avaliacao} resume os marcos mais relevantes para este
trabalho.

\begin{table}[htbp]
\centering
\caption{Marcos legais nacionais sobre avaliação de serviços públicos digitais}
\label{tab:marcos-legais-avaliacao}
\small
\begin{tabular}{p{3.2cm} p{1.4cm} p{8.2cm}}
\hline
\textbf{Instrumento} & \textbf{Ano} & \textbf{Principais disposições relevantes} \\
\hline
Lei n\textsuperscript{o} 12.527 (LAI) \cite{brasil2011lei12527} & 2011 &
Transparência ativa e passiva; base para acesso a informações públicas e dados
sobre desempenho de serviços. \\
Lei n\textsuperscript{o} 13.460 \cite{brasil2017lei13460} & 2017 &
Direitos do usuário, participação na avaliação dos serviços e uso dos resultados
para reorientação da prestação pública. \\
Decreto n\textsuperscript{o} 9.094 \cite{brasil2017decreto9094} & 2017 &
Simplificação do atendimento, Carta de Serviços ao Usuário e racionalização de
exigências feitas ao cidadão. \\
Lei n\textsuperscript{o} 14.129 \cite{brasil2021lei14129} & 2021 &
Governo digital, eficiência pública, desburocratização, canais digitais de
participação e melhoria da prestação estatal. \\
Decreto n\textsuperscript{o} 11.260 \cite{brasil2022decreto11260} & 2022 &
Elaboração e encaminhamento da Estratégia Nacional de Governo Digital, com
organização das diretrizes de transformação digital. \\
Portaria SGD/ME n\textsuperscript{o} 548 \cite{brasil2022portaria548} & 2022 &
Avaliação de satisfação dos usuários e padrões de qualidade para serviços
públicos digitais. \\
\hline
\end{tabular}

\footnotesize Fonte: elaborada pelos autores com base na legislação federal
brasileira.
\end{table}

Do ponto de vista técnico, relatórios produzidos no âmbito da Universidade de
Brasília, da Secretaria de Governo Digital e de instituições parceiras
apresentam modelos voltados à mensuração da qualidade de serviços públicos
digitais. A proposta Qualidade Gov.br organiza a avaliação em dimensões ligadas
à consistência digital e à experiência do usuário, reconhecendo que a qualidade
de um serviço digital depende tanto dos requisitos técnicos quanto da jornada do
cidadão \cite{pedrosa2019qualidadeGovbr}. O BRASP amplia essa discussão ao
considerar serviços presenciais, digitizados e parcialmente digitizados
\cite{silva2019brasp}.

O avanço desses modelos demonstra que a avaliação pública precisa equilibrar
rigor metodológico e simplicidade operacional. Instrumentos longos e complexos
podem gerar dados ricos, mas tendem a aumentar o esforço do respondente. Por
outro lado, instrumentos muito curtos podem ampliar a adesão, mas produzir dados
insuficientes para orientar melhorias. Esse equilíbrio é especialmente relevante
para a proposta deste TCC, que busca um avaliador simples para o cidadão e útil
para a gestão.

\section{BRASP, API de Avaliação e Modelo Unificado}

O Modelo Brasileiro de Avaliação de Serviços Públicos (BRASP) foi desenvolvido
para construir e validar instrumentos de avaliação da qualidade adaptados à
realidade brasileira. O modelo considera três modalidades de prestação:
serviços digitizados, serviços parcialmente digitizados e serviços presenciais
\cite{silva2019brasp}. Essa diferenciação é relevante porque a experiência do
usuário muda de acordo com o canal. Um serviço presencial pode depender mais de
atendimento, espera e estrutura física; um serviço digital pode depender mais de
facilidade de uso, clareza, disponibilidade e segurança; e um serviço híbrido
pode combinar problemas de ambos os contextos.

Nos materiais da Enap, o modelo brasileiro de avaliação é apresentado como
instrumento para identificar fatores relacionados à satisfação do usuário,
localizar pontos críticos e gerar comparações capazes de acompanhar a evolução
dos serviços. Essa leitura é coerente com o BRASP e com a proposta de modelo
unificado, que organizam satisfação, qualidade percebida, qualidade técnica e
efetividade como dimensões complementares de análise
\cite{limaPereira2020modulo3,silva2019brasp,silva2022modeloUnificado}. A
pesquisa de satisfação aparece ancorada em uma pergunta geral sobre a
experiência com o serviço, enquanto a qualidade percebida é investigada por
atributos que ajudam a explicar essa experiência.

A API de avaliação do governo brasileiro segue uma lógica semelhante. Segundo o
relatório sobre os atributos da API, a satisfação é medida por uma estratégia
simples inspirada no CSAT, com resposta por opções associadas a estrelas, e a
qualidade percebida é explorada por atributos como clareza das informações,
custo ou esforço para obtenção do serviço, disponibilidade, eficácia, facilidade
de uso, qualidade dos canais de comunicação e tempo para obtenção do serviço
\cite{silva2024atributosApi}. Essa abordagem se aproxima da proposta deste
trabalho, pois combina avaliação rápida com categorização do motivo da nota.

A proposta de modelo unificado reforça essa direção ao articular quatro eixos:
qualidade técnica, qualidade percebida, satisfação do usuário e efetividade
\cite{silva2022modeloUnificado}. Para um avaliador de serviços públicos, esses
eixos podem orientar tanto a coleta quanto a visualização de dados. A satisfação
indica o julgamento geral do cidadão; a qualidade percebida explica a experiência
por atributos; a qualidade técnica pode apoiar análises de disponibilidade,
integração e consistência; e a efetividade aproxima a avaliação da pergunta
central: o serviço resolveu a necessidade do usuário?

\section{Usabilidade, Inclusão Digital e Métricas de Experiência}

A usabilidade em serviços públicos digitais possui impacto mais amplo do que em
produtos privados. Quando uma interface privada é difícil, o usuário pode trocar
de fornecedor; quando um serviço público digital é difícil, o cidadão pode perder
prazo, deixar de acessar um benefício, não concluir uma obrigação ou depender de
atendimento presencial. Por isso, os princípios de usabilidade precisam ser
interpretados como parte do acesso efetivo ao serviço. As heurísticas de
\citeonline{nielsen1994heuristic}, como visibilidade do status do sistema,
correspondência entre linguagem da interface e linguagem do usuário, prevenção de
erros e ajuda contextual, são diretamente relacionadas aos problemas que o
avaliador busca observar durante a navegação.

A representatividade também é um ponto relevante. Serviços públicos atendem
usuários com diferentes níveis de letramento digital, conectividade, idade,
familiaridade com formulários e capacidade de compreender linguagem
institucional. Os dados do PPA Participativo 2024-2027 indicam que participação
digital em larga escala pode conviver com barreiras de uso e de retorno ao
cidadão \cite{brasilMgi2024ppaParticipativo}. Da mesma forma, pesquisas sobre
acesso às tecnologias de informação no Brasil mostram que o uso de serviços
digitais ainda depende de condições desiguais de acesso, dispositivos e
habilidades \cite{cetic2025ticDomicilios2024}. Para este TCC, isso reforça a
necessidade de uma coleta de feedback simples, opcional e de baixo esforço.

Métricas como NPS, CSAT e CES ajudam a decompor a experiência. O NPS popularizou
a pergunta de recomendação em escala de 0 a 10 \cite{reichheld2003nps}; em
serviços públicos, sua leitura precisa ser adaptada, pois o cidadão nem sempre
escolhe usar ou recomendar um serviço estatal. O CSAT registra satisfação
imediata com uma interação específica, normalmente por uma escala curta. O CES
mede o esforço percebido para concluir uma tarefa, aspecto especialmente
importante em serviços públicos com muitas etapas ou linguagem burocrática.
Coletadas no momento correto, essas métricas tendem a preservar melhor a
percepção do usuário do que avaliações muito afastadas da experiência, que podem
ser afetadas pelo peso do pico emocional e do desfecho da jornada
\cite{kahneman1993peakEnd}.

No protótipo, essa discussão aparece na combinação entre avaliação por estrelas,
atributos do problema e NPS. A estrela registra a percepção geral; os atributos
identificam se a dificuldade esteve ligada a clareza, tempo, esforço,
usabilidade, atendimento ou eficácia; e o NPS oferece um indicador comparável de
percepção global. A proposta não é transformar métricas privadas em objetivo
público, mas usar instrumentos simples para gerar dados interpretáveis pelo
gestor.

\section{Experiência do Usuário e Avaliações Inspiradas em Plataformas Digitais}

Plataformas digitais privadas popularizaram formas de avaliação de baixo esforço:
nota de uma a cinco estrelas, comentários curtos, seleção de motivos e
visualização agregada das avaliações. A referência a aplicações como plataformas
de entrega deve ser entendida como inspiração de interação, e não como modelo de
governança. Em uma plataforma privada, a avaliação frequentemente orienta
reputação, consumo e concorrência. Em serviços públicos, ela deve orientar
direitos, transparência, participação social e melhoria da entrega estatal.
Mecanismos como o NPS, proposto originalmente como pergunta de lealdade e
recomendação, também influenciaram práticas de avaliação de experiência por sua
simplicidade de aplicação \cite{reichheld2003nps}.

Mesmo com essa diferença, a familiaridade do usuário com mecanismos simples de
avaliação pode ser aproveitada. Um avaliador de serviços públicos pode reduzir
barreiras de participação se oferecer uma interface direta, com poucos passos e
linguagem compreensível. A nota geral permite captar a percepção imediata; os
atributos permitem interpretar os motivos; e comentários opcionais podem revelar
problemas não previstos pelo formulário. Essa escolha está alinhada ao design
centrado no usuário, que orienta o desenvolvimento de sistemas interativos a
partir das necessidades, capacidades e limitações das pessoas que irão utilizá-los
\cite{iso92412102019}.

Para que a avaliação seja útil ao gestor, o sistema também precisa organizar os
dados coletados. Isso envolve agrupar avaliações por serviço, período, órgão,
canal e atributo; identificar itens recorrentes; apresentar médias e distribuições;
e permitir comparação temporal. A literatura de melhoria de serviços reforça que
as avaliações devem alimentar ciclos de melhoria contínua, funcionando como linha
de base para decisões e como indicador para acompanhar a evolução do serviço
\cite{limaPereira2020modulo4,govuk2021usingData,govbrFerramentaAvaliacao}.

\section{Métricas de Navegação e Sinais de Dificuldade}

Além do feedback declarado pelo usuário, serviços digitais podem ser avaliados
por sinais observados durante a jornada. A taxa de conclusão, por exemplo, mede a
proporção de transações iniciadas que chegam ao ponto definido como conclusão do
serviço; sua leitura depende de identificar início, etapas intermediárias,
abandono e término do fluxo \cite{govuk2021completionRate}. Métricas de tempo,
cliques, páginas visitadas e retorno a etapas anteriores também podem indicar
dificuldade quando comparadas a faixas esperadas.

O uso dessas métricas aproxima a avaliação de serviços públicos da avaliação de
sistemas de informação. O modelo de sucesso de sistemas de informação de
\citeonline{deloneMclean2003} relaciona qualidade do sistema, qualidade da
informação, qualidade do serviço, uso, satisfação do usuário e benefícios
líquidos. Para a solução proposta, cliques e tempo de navegação funcionam como
indícios de uso e esforço; estrelas, atributos e NPS funcionam como percepção
declarada; e o painel administrativo busca transformar esses registros em
benefícios para a gestão.

O registro de eventos de navegação exige cuidado metodológico e jurídico. A
coleta deve ser limitada ao necessário para a finalidade de melhoria do serviço,
preferencialmente com sessão anônima ou identificador técnico sem exposição de
dados pessoais. A Lei Geral de Proteção de Dados Pessoais estabelece princípios
para o tratamento de dados, inclusive em meios digitais, com objetivo de proteger
liberdade, privacidade e desenvolvimento da personalidade da pessoa natural
\cite{brasil2018lgpd}. Por isso, a proposta deste TCC considera o registro de
cliques e tempo como mecanismo de análise agregada, não como instrumento de
vigilância individual.

\section{Feedback Contextual e Acionamento por Sinais de Uso}
\label{sec:feedback-contextual}

O feedback contextual diferencia-se do feedback episódico porque é solicitado no
momento em que a dificuldade ocorre ou logo após uma etapa relevante da jornada.
Em vez de depender de uma pesquisa enviada depois do uso ou de um formulário
isolado ao final do atendimento, a coleta contextual aproxima a pergunta da
experiência concreta. Essa proximidade reduz a perda de detalhe e evita que a
avaliação seja dominada apenas pela lembrança do desfecho da jornada
\cite{kahneman1993peakEnd}.

A contextualidade também envolve especificidade. Um comentário recebido após o
usuário permanecer tempo excessivo em uma etapa, clicar repetidamente em um
elemento ou retornar várias vezes ao mesmo caminho tende a indicar melhor o ponto
de atrito do que uma nota genérica posterior. Assim, o feedback deixa de ser
apenas um indicador de satisfação e passa a funcionar como insumo de melhoria de
produto: o gestor não recebe apenas a informação de que o serviço foi ruim, mas
também indícios sobre onde a dificuldade apareceu e qual atributo foi associado
ao problema.

No protótipo deste TCC, os gatilhos comportamentais foram reduzidos a regras
simples e auditáveis: excesso de cliques, tempo acima da faixa esperada e
continuidade da dificuldade após redirecionamento. Esses sinais não substituem a
percepção do usuário, mas indicam quando faz sentido solicitar uma avaliação de
baixo esforço. A escolha por regras explícitas preserva a rastreabilidade do MVP
e evita que a solução dependa de modelos estatísticos complexos antes de haver
dados reais suficientes para calibração.

Esse tipo de coleta também precisa respeitar a carga cognitiva do cidadão.
Questionários longos ou tecnicamente ambíguos aumentam esforço, erro e abandono.
A Teoria da Carga Cognitiva indica que a memória de trabalho humana possui
capacidade limitada; portanto, interfaces de feedback devem usar linguagem
direta, poucas perguntas, opções de resposta claras e campos textuais opcionais
\cite{sweller1988cognitiveLoad}. Por isso, a proposta adota estrelas, atributos
pré-definidos, NPS e comentário opcional como combinação mínima para gerar dados
úteis sem transformar a avaliação em nova barreira.

\section{Síntese do Referencial}

O referencial teórico indica que a proposta de um avaliador de serviços públicos
deve considerar quatro princípios. O primeiro é a centralidade do cidadão: a
avaliação precisa capturar a experiência real do usuário, e não apenas métricas
internas da administração. O segundo é a decomposição da avaliação em atributos:
uma nota geral é útil, mas precisa ser acompanhada de critérios que expliquem o
motivo da satisfação ou insatisfação. O terceiro é a orientação para melhoria:
os dados coletados devem ser estruturados para apoiar decisões de gestão e
acompanhar a evolução dos serviços. O quarto é a contextualidade: sempre que
possível, a avaliação deve ser solicitada próxima ao ponto de dificuldade, para
que o dado reflita melhor a experiência concreta.

Com base nesses princípios, este TCC propõe um avaliador inspirado em interfaces
simples de plataformas digitais, mas fundamentado em modelos de avaliação de
serviços públicos, usabilidade, design centrado no usuário e sucesso de sistemas
de informação. A solução deve permitir que o usuário avalie um serviço de forma
rápida e compreensível, enquanto eventos de navegação e resultados agregados
oferecem ao gestor indícios sobre quais aspectos do serviço precisam de
priorização e aperfeiçoamento.
